\documentclass[11pt]{article}
\usepackage[a4paper,margin=0.82in]{geometry}
\usepackage[T1]{fontenc}
\usepackage[utf8]{inputenc}
\usepackage{lmodern,amsmath,amssymb,amsthm,microtype,booktabs,tabularx,xcolor}
\usepackage[numbers]{natbib}
\usepackage[colorlinks=true,linkcolor=blue!60!black,citecolor=blue!60!black]{hyperref}
\setlength{\emergencystretch}{4em}
\newtheorem{theorem}{Theorem}
\newtheorem{remark}{Remark}
\title{\textbf{The Two Pointwise Fixed-Atom Parents: A Source-Locked Interface Comparison}}
\author{Liang Wang\\Huazhong University of Science and Technology}
\date{July 2026}
\begin{document}
\maketitle
\begin{abstract}
The two O161 parents have the same six-axis target signature, but different literal normalizations (q/T versus q/N).  After explicitly locking the common summand, specialization N=T proves direct twist implies the bad-endpoint target.  The reverse implication is not established. The narrower bad-endpoint parent is selected first.
The classification is \texttt{ROUTE\_\allowbreak{}DECISION\_\allowbreak{}L1}; no L2 arithmetic
progress or endpoint credit is claimed.
\end{abstract}

\section{Frozen target}
The selected route is \texttt{O161.bad\_\allowbreak{}endpoint\_\allowbreak{}pointwise\_\allowbreak{}fixed\_\allowbreak{}atom} \citep{TPC182}.  The required
signature is actual fixed-$h_0$ packet, named fixed atom, deterministic every
prefix, deterministic every scale, fixed-$X$ power at that atom, and actual
active support.  Throughout $h_0=2$ is only a source-backed data fact.

\section{Exact result}
\begin{theorem}
Source-locked specialization theorem: on the identical actual-core summand, with uniform constants, exponent, atom, scale range, and endpoint range, the all-prefix direct estimate with normalization q/N specializes separately at N=T for every bad endpoint T.  Thus direct implies bad endpoint.  The reverse implication is not established.
\end{theorem}
\begin{proof}
Fix any T in the bad-endpoint set.  The direct hypothesis is uniform for every admissible prefix N, so take N=T.  The summand, atom, constant and exponent are unchanged, and q/N=q/T.  This proves the bad-endpoint bound for each such T.  The source graph contains no reverse edge and the bad hypothesis has the smaller endpoint domain; we therefore record the reverse implication as not established, not false.
\end{proof}

The smallest missing literal theorem is
\texttt{UNIFORM\_\allowbreak{}NAMED\_\allowbreak{}ATOM\_\allowbreak{}BAD\_\allowbreak{}ENDPOINT\_\allowbreak{}PREFIX\_\allowbreak{}POWER\_\allowbreak{}SAVING}.

\section{Scoped stop and claim firewall}
No new scoped stop is declared in this paper.

The result is L0/L1 only.  Phase averages and Lebesgue-almost-every phase are
not named atoms.  Fixed $h_0=2$ is not decay.  Archive addressing is not
canonical physical representation, and empty-domain truth is not actual
active support.  The named-atom endpoint charge is zero, so the strict
$1/400$ budget is unpaid.  We claim no program-positive L2, prime-pair lower
bound, or twin-prime theorem.

\section{Reproducibility}
The adjacent JSON payload records the full six-axis signature, source lock,
typed result, exact scoped-stop cell, endpoint ledger, and claim firewall.
The audit artifact contains eight true mutation-rejection assertions.
\bibliographystyle{plainnat}
\bibliography{references}
\end{document}
